\documentclass[10pt,a4paper]{amsart}
\usepackage[margin=19mm]{geometry}
\usepackage{amsmath,amssymb,mathtools}
\usepackage[scr=rsfs,cal=euler]{mathalfa}
\usepackage{xfrac}
\usepackage[unicode,hidelinks,bookmarks=false]{hyperref}
\newcommand{\ie}{i.e.\ }
\newcommand{\cf}{cf.\ }
\newcommand{\resp}{respectively\ }
\newcommand{\Subsection}{\subsection}
\providecommand{\nobreakdash}{\nobreak\hbox{-}}
\setlength{\emergencystretch}{2em}
\multlinegap0pt
\usepackage{bm}
\usepackage{bbm}
\usepackage{dsfont}
\usepackage{xspace}
\usepackage{cancel}
\usepackage{mathtools}
\usepackage{amsthm}
\makeatletter
\@ifundefined{theorem}{\newtheorem{theorem}{Theorem}}{}
\@ifundefined{lemma}{\newtheorem{lemma}[theorem]{Lemma}}{}
\makeatother
\title[Long-memory Markov chains with power-law intensities]{Long-memory Markov chains with power-law intensities}
\author{Kyungsub Lee}
\date{Source: arXiv:2607.20838v1 (CC BY 4.0)}
\begin{document}
\maketitle

\noindent\emph{Source and licence.} Long-memory Markov chains with power-law intensities. Version arXiv:2607.20838v1 (CC BY 4.0), distributed under the Creative Commons Attribution 4.0 International licence, as recorded in the arXiv licence metadata for this exact version and shown on the abstract page https://arxiv.org/abs/2607.20838v1. Full rights notice: licenses/arxiv-2607.20838-CC-BY-4.0.txt.\par\smallskip

\noindent\textit{Editorial scope.} Counted unit: the Lemma whose source label is \texttt{Lemma:derivatives} in the article (source0.tex lines 1076--1090, source label \texttt{Lemma:derivatives}), delivered together with the article's own complete original proof of it (source0.tex lines 1091--1190, one \texttt{proof} environment of 100 lines). No mathematics is written, repaired or re-derived: statement and proof are the article's own text line for line. Required context retained uncounted: Eq:1\_0 (lines 1066--1069); Eq:2\_0 (lines 1071--1074). Every definition, display and label of those lines is retained unchanged. Omitted from this package and not redistributed: 1--1065, 1070--1070, 1075--1075, 1191--1206. Nothing inside a retained excerpt is omitted or altered: every retained body is delivered exactly as the article's lines are written, so no omission receipt of any kind accompanies this package. Citations: the retained excerpts carry no citation command at all, so no bibliography entry of the article is transcribed and no reference is redistributed. Reference adaptation: none was needed and none was made; every reference inside the delivered unit resolves inside it, so no unresolved reference or citation is produced. Printed slips kept literal: no printed word, symbol, hypothesis or number of the counted statement or of its proof was altered. Layout and class: the article's own class is \texttt{article} with packages including natbib, graphicx, multirow, amsmath, amssymb, amsfonts, amsthm, mathrsfs, appendix, xcolor, textcomp, manyfoot, booktabs, algorithm; the wrapper is the lane's pinned \texttt{amsart} editorial wrapper at 10pt on A4, which restates the theorem-like environments the retained text needs and adds the article's own macro definitions that the retained text needs (none), with extra wrapper packages (bm, bbm, dsfont, xspace, cancel, mathtools). The delivered numbering is the unit's own. Assets: no retained span contains a figure environment, a graphics-inclusion command, a PDF-page inclusion, a table or a TikZ picture; no image file is copied into this package, the package declares no asset dependency and no image file is redistributed with it. Licence: version arXiv:2607.20838v1 of this article is distributed under the Creative Commons Attribution 4.0 International licence, as recorded in the arXiv licence metadata for this exact version and shown on the abstract page https://arxiv.org/abs/2607.20838v1. A full rights notice is delivered at licenses/arxiv-2607.20838-CC-BY-4.0.txt. Characters: every retained excerpt is pure ASCII, so the delivered engine, XeLaTeX with its default Latin Modern text font, reports no missing character.
\begin{equation}
	\left. X_1 \right|_{\mathbf{0}} = x + \xi,  \qquad
	\left. D_1 \right|_{\mathbf{0}} = \frac{\gamma x + c \xi}{\gamma c}, \qquad \left. c_1 \right|_{\mathbf{0}} = \frac{\gamma c(x + \xi)}{\gamma x + c \xi}, \label{Eq:1_0}
\end{equation}

\begin{equation}
\left. X_2 \right|_{\mathbf{0}} = x + 2 \xi,\qquad \left. D_2 \right|_{\mathbf{0}} = \frac{\gamma x + 2 c \xi}{\gamma c}, \quad \left. c_2 \right|_{\mathbf{0}} = \frac{\gamma c(x + 2\xi)}{\gamma x + 2 c \xi}. 
\label{Eq:2_0}
\end{equation}

\begin{lemma}\label{Lemma:derivatives}
The evaluated derivatives of of $X_2$ and $c_2$ with respect to $u_i$ at $u_1 = u_2 =0$ are
\begin{align}
&\left.\frac{\partial X_2}{\partial u_1}\right|_{\mathbf{0}} = -\frac{px}{c}, \quad\quad\quad\quad \left.\frac{\partial X_2}{\partial u_2}\right|_{\mathbf{0}} =  -\frac{p(\gamma x+c\xi)}{\gamma c}, \label{Eq:eval_1}\\
&\left.\frac{\partial c_2}{\partial u_1}\right|_{\mathbf{0}}
= \frac{\gamma x\big(\gamma x + 2p\xi(\gamma-c) + 2\gamma\xi\big)}{(2c\xi+\gamma x)^2} , \quad\quad\quad\quad
\left.\frac{\partial c_2}{\partial u_2}\right|_{\mathbf{0}}
= (\gamma x + c\xi) \frac{B}{M}	 \label{Eq:eval_2}
\end{align}
where
\begin{align*}
&B = \gamma x^2 + (p+2)\gamma \xi x + (1-p)c\xi x + 2c\xi^2, \\
&M = (x+\xi) (\gamma x + 2 c\xi)^2.
\end{align*}
\end{lemma}

\begin{proof}
	We begin with Eq.~\eqref{Eq:eval_1}.
	First, note that
	$$ \frac{\partial X_2}{\partial c_1} = p X_1 r_2^{-p-1} u_2 c_1^{-2}.$$
	Hence,
	\begin{equation}
	\left.\frac{\partial X_2}{\partial u_1}\right|_{\mathbf{0}}
	=
	\left.\frac{\partial X_2}{\partial X_1}\right|_{\mathbf{0}}\left.\frac{\partial X_1}{\partial u_1}\right|_{\mathbf{0}}
	+
	\left.\frac{\partial X_2}{\partial c_1}\right|_{\mathbf{0}}\left.\frac{\partial c_1}{\partial u_1}\right|_{\mathbf{0}}
	=
	1\cdot\left(-\frac{p x}{c}\right)+0\cdot(\cdots)
	=
	-\frac{p x}{c}, \label{Eq:X2_u1_0}
	\end{equation}
	Similarly,
	\begin{equation}
	\frac{\partial X_2}{\partial u_2} = 
	X_1\cdot(-p) r_2^{-(p+1)}\cdot\frac{1}{c_1}
	\quad\Rightarrow\quad
	\left.\frac{\partial X_2}{\partial u_2}\right|_{\mathbf{0}}
	= -\frac{p X_1}{c_1}\bigg|_{\mathbf{0}}
	= -\frac{p(\gamma x+c\xi)}{\gamma c}. \label{Eq:X2_u2_0}
	\end{equation}

	For the first expression in Eq.~\eqref{Eq:eval_2}, observe that
	\begin{equation}
	\left.\frac{\partial c_2}{\partial u_1}\right|_{\mathbf{\mathbf{0}}}
	= \left.\frac{\partial c_2}{\partial X_1}\right|_{\mathbf{0}}\left.\frac{\partial X_1}{\partial u_1}\right|_{\mathbf{0}} + 
	\left.\frac{\partial c_2}{\partial c_1}\right|_{\mathbf{0}}\left.\frac{\partial c_1}{\partial u_1}\right|_{\mathbf{0}}. \label{Eq:dc2_du1}
	\end{equation}
	From Eqs.~\eqref{Eq:1_0}, \eqref{Eq:2_0}, and
	\begin{align*}
	&\frac{\partial {X_2}}{\partial X_1} = r_2^{-p} \quad\Rightarrow\quad  \left. \frac{\partial {X_2}}{\partial X_1} \right|_{\mathbf{0}} = 1 \\ 
	& \frac{\partial{D_2}}{\partial X_1} =  - D_1 (p+1) r_2^{-(p+2)}\frac{\partial r_2}{\partial X_1} =  (p+1) r_2^{-(p+2)}\frac{u_2 D_1^2}{X_1^2} \quad\Rightarrow\quad  \left. \frac{\partial {D_2}}{\partial X_1} \right|_{\mathbf{0}} = 0,
	\end{align*}
	we obtain
	\begin{align}
		\left.\frac{\partial c_2}{\partial X_1}\right|_{\mathbf{0}} = \left. \frac{(\partial_{X_1}X_2) D_2 - X_2 (\partial_{X_1}D_2)}{D_2^2} \right|_{\mathbf{0}} = \left. \frac{1}{D_2} \right|_{\mathbf{0}} =  \frac{\gamma c}{\gamma x + 2c\xi }. \label{Eq:dc2_dX1}
	\end{align} 
	Also,
	\begin{equation}
	\left. \frac{\partial X_1}{\partial u_1} \right|_{\mathbf{{0}}} = \left. - x p r_{1}^{-p-1}\frac{1}{c} \right|_{\mathbf{{0}}}= -\frac{px}{c}. \label{Eq:dX_1_du_1_0}
	\end{equation}
	
	Next, since
	\begin{align*}
		\frac{\partial c_2}{\partial c_1} = \frac{\partial}{\partial c_1} \left(\frac{X_2}{D_2}\right) = \frac{(\partial_{c_1}X_2) D_2 - X_2 (\partial_{c_1}D_2)}{D_2^2} ,
	\end{align*}
	and
	$$ \left. (\partial_{c_1}X_2) \right|_{\mathbf{0}} = \left. pX_1 r_2^{p-1} u_2 \frac{1}{c_1^{2}} \right|_{\mathbf{0}} = 0, \qquad  \left. (\partial_{c_1}D_2) \right|_{\mathbf{0}} = \left. - \frac{X_1}{c_1^2} \right|_{0} = - \frac{(\gamma x + c \xi)^2}{(c\gamma)^2 (x + \xi)},
	$$
	by Eq.~\eqref{Eq:2_0},
	we have
	\begin{equation}
	\left.\frac{\partial c_2}{\partial c_1}\right|_{\mathbf{0}}
	= \frac{(x+2\xi)(\gamma x + c\xi)^2}{(x+\xi)(\gamma x + 2c\xi)^2 }. \label{Eq:dc2_dc1}
	\end{equation}
	From Eqs.~\eqref{Eq:1_0} and \eqref{Eq:dX_1_du_1_0} and using
	$$\left. \partial_{u_1} D_1 \right|_{\mathbf{0}} = - \frac{(p+1)x}{c^2},$$
	we have
	$$
	\left. D_1 (\partial_{u_1} X_1) - X_1 (\partial_{u_1} D_1) \right|_{\mathbf{0}}
	=  -\frac{px}{c} \frac{x \gamma + \xi c}{c \gamma} + \frac{(p+1)x}{c^2}(x + \xi )
	$$
	and
	\begin{equation}
	\left.\frac{\partial c_1}{\partial u_1}\right|_{\mathbf{0}}
	=\left. \frac{D_1 (\partial_{u_1} X_1) - X_1 (\partial_{u_1} D_1) }{D_1^2} \right|_{0}=  \frac{x\gamma\big(\gamma x + (p+1)\gamma\xi - p c\xi\big)}{(\gamma x + c \xi)^2}. \label{Eq:dc1_du1}
	\end{equation}
	Substituting Eqs.~\eqref{Eq:dc2_dX1},\eqref{Eq:dX_1_du_1_0},\eqref{Eq:dc2_dc1} and \eqref{Eq:dc1_du1} into Eq.~\eqref{Eq:dc2_du1} yields the desired expression for $\partial c_2 / \partial u_1$.

	For the second equality in Eq.~\eqref{Eq:eval_2}, note that
	$$
	\frac{\partial c_2}{\partial u_2}
	=\frac{(\partial_{u_2}X_2) D_2 - X_2 (\partial_{u_2}D_2)}{D_2^2}.
	$$
	From Eq.~\eqref{Eq:X2_u2_0} and
	$$
	\partial_{u_2} D_2 = \frac{X_1}{c_1} (-(p+1)) r_2^{-(p+2)}\frac{1}{c_1}
	\quad\Rightarrow\quad
	\left.\partial_{u_2}D_2\right|_{\mathbf{0}} = \left. -\frac{(p+1) X_1}{c_1^2} \right|_{\mathbf{0}},
	$$
	we find
	\begin{align*}
	 \left. (\partial_{u_2}X_2)D_2 - X_2(\partial_{u_2}D_2) \right|_{\mathbf{0}}
	&= \left. \left(-\frac{p X_1}{c_1}\right)D_2 - (X_1+\xi)\left(-\frac{(p+1) X_1}{c_1^2}\right) \right|_{\mathbf{0}} \\
	&= \left. \frac{X_1}{c_1}\left( (p+1) \frac{X_1 + \xi}{c_1} - p D_2 \right) \right|_{\mathbf{0}} \\
	&=  \frac{\gamma x + c\xi}{c \gamma } 
	\left( (p+1)\frac{(x+2\xi)(\gamma x + c\xi)}{c\gamma(x+\xi)} -p\frac{\gamma x+2c\xi}{c \gamma}\right).
	\end{align*}
	Thus,
	$$
	\left.\frac{\partial c_2}{\partial u_2}\right|_{\mathbf{0}}
	= \frac{(\gamma x + c\xi)}{(x+\xi)(\gamma x + 2c\xi)^2}
	\Big( \gamma x^2 + ((p+2)\gamma + (1-p)c ) \xi x + 2c\xi^2 \Big),
	$$
	which completes the proof.
\end{proof}

\end{document}
